%!TEX root = /Users/markelikalderon/Documents/Git/perceptual_self-consciousness/perceptual_self-consciousness.tex
\chapter{Return, and So Coming, at Last, to Rest} % (fold)
\label{cha:return}

\section{Introduction} % (fold)
\label{introduction}

In this, the final chapter, drawing upon the insights gleaned from our respected predecessors, I shall outline an account of perceptual self-consciousness.

on non-positional self-awarenesss and relate that account to soem of the many issues that arise in contemporary discussions of self-consciousness.

% section introduction (end)

\section{Some Preliminary Obstacles} % (fold)
\label{some_preliminary_obstacles}

Notoriously, \citet[324]{Armstrong:1968ns} makes the following bold denial: ``A mental state cannot be aware of itself, any more than a man can eat himself.'' (Armstrong, like many philosophers of mind, speaks indifferently of mental states or events. On the perils of this indifference see \citealt{Steward:1997aa}.) A man cannot eat himself, or, at least, the whole of himself. But is the awareness afforded by a mental state really analogous? What is the basis of this forthright denial?

\citet[324]{Armstrong:1968ns} insists that, in the case of perception, ``we must distinguish between the perception\ldots from the thing perceived.'' So the thing perceived is other than the perception of it. So like Sartre, Armstrong maintains that the intentional object of perception is transcendent in being other than the perception directed upon it.

In denying that a mental state can be aware of itself, Armstrong does not thereby deny that we may be directly aware of our own mental states. Armstrong merely denies that the source of this awareness is immanent to the mental state of which we are aware. Rather, the source of the direct awareness of our mental states is introspection, a distinct higher-order mental state. Introspection, for Armstrong, is quasi-perceptual:
\begin{quote}
	In introspection\ldots we become aware of current happenings in our own mind. Introspection is therefore a mental event having as its (intentional) object other mental happenings that form part of the same mind. Nevertheless, introspection may properly be compare to sense-perception, and Kant's description of introspection as `inner sense' is perfectly justified. \citep[323]{Armstrong:1968ns}
\end{quote}
Given its root, Armstrong is right to confine talk of intro\underline{spection} to a quasi-perceptual awareness of mental happenings. (Not all philosophers have been so discriminating; on introspection see \citealt{Lyons:1986fe} and \citealt{Spener:2024zh}). It is the differences between perception and intropsection that warrants the qualification. Perception depends upon sensory organs that are subject to voluntary control \citep[324]{Armstrong:1968ns}. Let the latter claim be liberally undertood. We cannot cock our ears like dogs, even if we can voluntarily reposition ourselves to better listen. Intropsection, by contrast, has no dedicated organ subject to voluntary control. In this regard, introspection is less like the external senses than forms of bodily awareness such as propioception and kinaesthesis \citep[324]{Armstrong:1968ns}. This is why Armstrong cites with approval Kant's description of introspection as an ``inner sense''. But the tradition is much older, and Armstrong is the unwitting heir to Ibn Sina. As we shall see, there are other elements of his view with ancient precedent.

Just as perception is distinguished from the thing perceived, \citet[324]{Armstrong:1968ns} distinguishes introspection and the thing introspected: ``it is impossible that that the introspecting and the thing introspected should be one and the same mental state.'' Despite both being mental states, the thing introspected is other than the introspection. The intentional object of introspection is transcendent in being other than the introspecting directed upon it.

Like Socrates (not only in the \emph{Charmides} but also in the \emph{Gorgias}, 449d, 454a, and the \emph{Symposium}, 199e), Armstrong maintains that the intentional object of a mental state is distinct from that state. Moreover, like Socrates (\textsc{The Argument from Relatives}, \emph{Charmides} 168b2–169c), Armstrong is moved by an apparent tension between intentionality and reflexivity. It is because an intentional object is transcendent that a mental state cannot be aware of itself, any more than a man can eat himself. 

Rejecting reflexivity is not to reject self-knowledge or perceptual self-conscious\-ness but only a conception of what these amount to. These are genuine mental phenomena, but, according to Armstrong, they are, being intentional, misunderstood as reflexive. Rather, they are introspective.

Armstrong may be right about the tension between intentionality and reflexivity, but Armstrong is wrong if he thinks that sufficient grounds to deny that a mental state may be aware of itself. Recall, there were two ways out of the impasse posed by \textsc{The Argument from Relatives} (chapter~\ref{sec:coda}). Given the \emph{aporia}, one could either hang onto intentionality and reject reflexivity or hang onto to reflexivity and reject intentionality. If a mental state, in being aware of itself, does not thereby take itself as an intentional object, then one can accept the transcendent nature of intentionality consistent with reflexivity. Reflexive self-awareness is coherent only if it is  non-positional. 

% section some_preliminary_obstacles (end)

% chapter return (end)